\subsection{The unique action on four orthogonal blocks}

\begin{statement}{D23.P21}{Proved}{A central phase and its actual two-cycle normalization}
\end{statement}
Restrict the first two braid generators of D23.P17 to the
union of orthogonal blocks $7,8,10,11$, of total degree
448. The resulting $B_3$ image has order 11664 and maps
onto the degree-36 detector image of order 3888. Its kernel
is the central subgroup of order three generated by
\[
                         z=(AB)^3.
\]
The pure-pair image $P=\langle A^2,B^2\rangle$ has order
5832, and $P'$ has order 1944. The same central kernel is
contained in $P'$.

Each of the six actual first-pair values determined by
D23.H1 has exactly three lifts in $P'$. Exactly one lift
satisfies the two-cycle. Consequently the entire first-pair
value on these 448 points is determined by $(\rho,a_3)$.

For $\rho=0$, use the H3 candidate indices and block order
$(7,8,10,11)$. The unmodified first generator has already
reduced the possibilities to $(i,i,i,j)$ in D23.P20.
Covariance of the second generator further gives
\[
\begin{array}{c|c}
a_3&\text{allowed pairs }(i,j)\\ \hline
0&(0,0),(1,1),(2,2)\\
1&(0,2),(1,1),(2,0).
\end{array}                                                \tag{D23.P21.1}
\]
Combining this with the independently computed block-11
condition in D23.H4 fixes the action on all four blocks:
\[
\begin{array}{c|c}
a_3& (i_7,i_8,i_{10},i_{11})\\ \hline
0&(0,0,0,0)\\
1&(0,0,0,2).
\end{array}                                                \tag{D23.P21.2}
\]
In particular the normalized mixed actual full-GT state,
if it exists, has the unique indicated nontrivial action
on these four blocks. Existence is not asserted.

\begin{proof}
The detector map on the 432 indicated points is an exact
$B_3$-equivariant map to the degree-36 action. The union
of the four orthogonal blocks is invariant under $A,B$;
the remaining sixteen points cause no ambiguity in the
map between the permutation image groups. Exact
stabilizer-chain calculations give the four orders
\[
 |\langle A,B\rangle|=11664,\quad
 |\langle A_{36},B_{36}\rangle|=3888,\quad
 |P|=5832,\quad |P'|=1944.
\]
The images of $P$ and $P'$ on the detector have orders
1944 and 648 by D23.H1. Hence each of the three maps
has kernel of order three. Direct evaluation shows that
$z=(AB)^3$ is nonidentity, has order three, projects to
the identity on the detector, commutes with $A,B$, and
belongs to $P'$. It therefore generates every one of
these kernels. The certificate verifies all of these
facts, including membership in $P'$.

Put $s=ABA$. The braid relation gives
\[
        sA^2s^{-1}=B^2,\qquad sB^2s^{-1}=A^2.
\]
For every actual commutator associator $f$ satisfying
the two-cycle, its value $F=f(A^2,B^2)$ consequently
satisfies
\[
                         sFs^{-1}=F^{-1}.                 \tag{D23.P21.3}
\]
If one lift $F$ of a prescribed detector value satisfies
this equation, then replacing it by $Fz^k$ changes the
two sides by $z^k$ and $z^{-k}$, respectively. Since
$z$ is central of order three, at most one of the three
lifts can satisfy the equation. Enumerating the 1944
elements of $P'$ finds exactly three lifts of each of
the six H1 values, and exactly one of them satisfies
(D23.P21.3). All six selected values also satisfy the
cyclic norm condition under $AB$; this is a check of
consistency, not a sufficiency claim for actual GT.

Now form each of the nine permutations from D23.P20
on the 448-point set and test the exact covariance
\[
                     VBV^{-1}=F^{-1}BF,
\]
using the selected $F$ for its $(\rho,a_3)$ value.
This complete nine-element calculation yields exactly
(D23.P21.1). It uses the literal permutations on all
448 points, not only their orthogonal matrices or their
degree-36 quotients.

Finally D23.H4 proves that at $\rho=0$ the mixed value
$R=f(X_3,X_4)$ is the identity on block 11. Thus $V_{11}$
must commute with the local third braid generator. Its
unique H3 candidate is index zero for $a_3=0$, and
index two for $a_3=1$. Substitution in (D23.P21.1)
gives (D23.P21.2).
\end{proof}

\paragraph{Certificate and exact input locators.}
The full finite certificate is
\path{hurwitz_first_four_block_lifts.py}, with JSON and
NPZ records of the same stem. The NPZ contains the
448-point ordering, its two generators, both
$\rho=0$ first-pair values, and every permutation
passing the displayed covariance test.

The actual detector values come from D23.H1,
\path{work/records/hurwitz36_gt_stability.json}.
The exhaustive orthogonal lift lists come from D23.H3,
\path{work/notes/hurwitz_h1_orthogonal_lifts.tex} and its
code/records. The final block-11 condition is D23.H4,
\path{work/notes/hurwitz_block11_harmonic.tex}, with
\path{work/code/hurwitz_block11_harmonic.py} and its
JSON/NPZ records. These are necessary actual-source
constraints, not constructions of an associator from
the finite values.
